\section{Discussion}
\label{sec::chp6::discussion}

One of the primary strengths of \deeid is in its extension of the Fiat--Shamir identification protocol to accommodate multiple identity issuers. Unlike traditional systems that rely on a single trusted authority or on centralised certificate authorities, \deeid's use of blockchain as a DPKI enables a dynamic trust network where any organisation—be it a bank, government, or university—can issue verifiable credentials. This decentralisation eliminates single points of failure and reduces the risk of systemic compromise, as demonstrated by real-world incidents such as the SEC Twitter/X account hack via SIM swapping \cite{kan_hijack_2024} or the deepfake-enabled US\$25 million fraud \cite{milmo_uk_2024}. By replacing shared secrets with zero-knowledge proofs, \deeid ensures that no reusable information is exposed during verification, significantly reducing the attack surface for impersonation and replay attacks.

%The system's compatibility with existing web technologies is another critical advantage. By utilizing QR codes, WebSocket communication, and standard JavaScript or Python libraries, \deeid integrates seamlessly with current web infrastructure without requiring browser modifications or proprietary software. This deployability makes it practical for organizations to adopt \deeid without overhauling their existing systems, addressing a common barrier to the adoption of advanced authentication schemes \cite{bonneau_quest_2012}. Furthermore, the support for multiple authentication methods—ranging from explicit blockchain-based links to anonymous public-key generation—offers flexibility to cater to diverse use cases, from high-security scenarios requiring strong identification to privacy-focused applications needing unlinkable authentication.

\deeid has taken a privacy-by-design approach. By employing zero-knowledge proofs, the system ensures that verifiers receive only the information necessary for authentication/identification, protecting users from oversharing sensitive data. This aligns with growing regulatory demands for data minimisation, such as those outlined in the GDPR and similar frameworks. Additionally, the ability to support multiple devices and keys enhances usability, allowing users to manage their identities across various contexts without being locked into a single vendor or ecosystem, unlike federated identity systems.

%Human identification is an integral part of the security of the digital ecosystem. Though it is a missing element at the moment which is leading to account breaches, identity theft, and much more. In this chapter, we introduced the Fiat--Shamir scheme and have generalised it so that it could work on a larger scale with \textit{multiple} identity issuers. Our solution has been to use the blockchain as a decentralised PKI. Blockchain-based PKIs have shown promise in research and applications. Emergence of blockchain-based PKIs have come about mainly due to the weak-link security problem with respect to certificate authorities and the X.509 standard~\cite{noauthor_x.509_nodate} (other main PKI standard is OpenPGP~\cite{shaw_openpgp_nodate}). Consequently, this has led to new approaches such as the Google Certificate Transparency project~\cite{noauthor_certificatetr_nodate}. Blockchain PKI models such as~\cite{konoplev_blockchain_2018} attempt to overcome the weaknesses of current PKI standards.

% Essentially we have a p2p network and we have to represent the human identities on this network. How should the identity be represented? and what mechanisms are required to interact with the identity? Going back to our four rules: a) Capture: each identity must have a unique representation on the network. Each \deeid\ contract has a unique address. b) Recording: It must be machine and human readable. The \deeid\ contract address is unique and machine-readable. The user can also prove their identity and give their human readable names to the challenger. c) Security: can we protect from impersonation and other attacks? Best current method is using public-key cryptography and \textit{what we posses} rather than \textit{what we know} or \textit{what we are}. d) Verification: Using the mobile phone device and the Fiat--Shamir scheme we can verify identities easily.

%Similar to uPort we have represented the identity by a smart-contract which is uniquely identified in the network. Therefore, each identity has an interactive code on the network (managed by the identity holder themselves). So, here, identity is no longer a piece of string such as your full name, but a totally unique digital representation on a network. The benefit of this is that it is a single-source of truth that everyone can use and agree upon. So, your bank, hospital and school can all use the same identity. Using the same identity and authentication scheme can provide data sharing which in itself has huge benefits. We believe merging the Fiat--Shamir scheme with the existing blockchain identity methods is a novel method for an identification system.

%The benefit of using the Fiat--Shamir scheme over simply signing an identity string is the privacy and security it provides. Moreover it saves transactions and space on the network. However, it is still important to keep our 'full names', the basic string identity, because we humans still need to communicate with one another and be able to identify each other. Which is not possible if we are identified by a 512bit unique identity string. There have been studies such as \cite{bonneau_towards_2014}, \textit{'Towards reliable storage of 56-bit secrets in human memory'}, where the authors tried to see if we humans can remember long bits of string. The conclusion was that humans are able to learn cryptographic secrets (56-bit in this case), though it requires significant learning periods and other limitations such as recalling times.

%The issue of how do we trust the identity issuer falls upon the verifier. At the moment it is conceived to be a simple linear trace back of identities until the final link must somehow be verified. In our implementation we have a simple link to an official website or regulatory-body portal. However, there is considerable work to be done in improving this, e.g. one can quantify trust and provide recommendation to the verifier whether they are trust-worthy. A reputation-based network can provide some metric for verifiers to work from. However, one must consider the ethical implications of quantifying reputation, especially when we are considering real identities.

%\deeid predates Passkeys but offers similar advantages to it and WebAuthn when it comes to password-less authentication. Traditional authentication systems expose users to compromise when websites, third-party libraries, browsers, or operating systems contain vulnerabilities, as demonstrated by the 2018 British Airways breach where attackers compromised website scripts to steal customer data during entry, including CVV numbers that organizations typically never store~\cite{noauthor_british_2019, bishop_this_nodate}. In contrast, \deeid eliminates these input-capture vulnerabilities by using QR code scanning to initiate encrypted WebSocket connections that bypass browser-based data entry entirely, preventing scrapers and compromised web components from intercepting authentication credentials. The system employs public-key cryptography and digital signature schemes rather than passwords, providing four distinct authentication methods through public keys stored on users' identity smart contracts. This architecture enables users to manage their own security posture by adding or removing cryptographic keys as needed, offering both enhanced security and user autonomy compared to centralised password-based systems that remain vulnerable to website-level compromises.

When compared to existing authentication systems like WebAuthn and Passkeys, \deeid stands out for its focus on identification rather than solely authentication. While WebAuthn effectively addresses password-less authentication through public-key cryptography, it does not inherently support the complex, multi-issuer identity verification scenarios that \deeid targets. For instance, WebAuthn is primarily designed for single-site authentication, whereas \deeid enables cross-organisational identity verification through its blockchain-based trust network. This makes \deeid particularly suited for applications requiring robust identity verification, such as financial services, healthcare, or government services, where multiple trusted issuers are involved.

Traditional public key infrastructures (PKIs) rely on centralised certificate authorities, which introduce vulnerabilities such as single points of failure and potential coercion. In contrast, \deeid's blockchain-based DPKI distributes trust across a decentralised network, reducing reliance on any single entity. Compared to other blockchain-based identity systems like Namecoin \cite{noauthor_namecoin_nodate} or Certcoin \cite{fromknecht_certcoin_2014}, \deeid's integration of the Fiat--Shamir protocol with a multi-issuer framework provides a more flexible and scalable solution for real-world identity management. Additionally, \deeid's emphasis on zero-knowledge proofs distinguishes it from systems like Blockstack \cite{noauthor_blockstack_nodate}, which may not prioritise privacy to the same extent.


% In our evaluation we followed the framework described in \cite{bonneau_quest_2012} to do a comparative analysis of \deeid with existing authentication schemes. Schemes such as: passwords, password managers, federated systems, hardware, mobile phone based systems, biometrics and more. The \deeid protocol uses cryptographic keys that are unpractical for humans to read, copy and transfer themselves without the help of a machine. Moreover, the computation required for these keys is impossible for humans to do. Thus, the use of a machine such as a mobile phone is required. Consequently, our protocol blurs the lines between hardware tokens, phone-based systems, federated and password management systems as shown in Figure~\ref{fig:deeIDArchitecture}.

% It should be noted that the framework in Figure~\ref{fig:deeIDArchitecture} is designed for human practicality, incorporating factors such as ``nothing to carry'', ``physical-effortlessness'' and ``memory-wise-effortless''. Therefore, the protocols are highly human-centric at present. Nevertheless, as we progress into the domains of AI and IoT, we can improve upon these protocols and frameworks to accommodate other non-human entities as well. Whilst no other system can surpass passwords that are easily stored and retrieved in our brains, as the number of passwords to remember increases, we encounter difficulties. Thus, the benefits of alternative schemes, such as password managers, become more apparent.


Preventing SIM swap attack is another use case of \deeid. SIM swap attack is where an attacker uses social-engineering (operator error and weakness) to convince the user's mobile carrier to swap the SIM to the attacker's SIM. The consequence of this is that the attacker can hijack accounts that use two-factor authentication (2FA). Notable hacks include the hijacking of celebrity Instagram accounts~\cite{bell_instagram_nodate} and individuals losing millions of pounds of cryptocurrency as their exchange accounts were linked to their mobile phone number using 2FA. This is a growing threat that is leading to more organised crime~\cite{franceschi-bicchierai_sim_2018}. Preventive techniques include using a PIN or better 2FA~\cite{wire_protect_swimswap}. However, these are patchy and preventive measures; a standardised solution through the use of decentralised identity like that of \deeid\ across mobile network carriers would be a better solution. Explicitly linking the user's \deeid\ to their SIM and having a hardware linked crypto keys - meaning that it is physical and malicious attackers cannot carry out their attack without access to the physical elements. When the SIM is linked to the user's \deeid\ (do this when the user acquires the SIM), the attacker would have to identity themselves which is far more difficult that knowing the user's name, address and date of birth.

Mitigating identity theft presents a compelling use case for \deeid, as current credit verification systems enable attackers to obtain credit cards using only a victim's name, address, and date of birth. This vulnerability stems from credit agencies operating on a ``pull'' model, where third parties can access credit reports by providing basic personal information, rather than a ``push'' model requiring explicit authorization from the identity owner. Under a \deeid-based framework, credit agencies such as Experian~\cite{noauthor_experian_nodate}, Equifax~\cite{noauthor_equifax_nodate}, and TransUnion~\cite{noauthor_transunion_nodate} would link customer data to individual \deeid credentials, transforming credit inquiries into authorization requests that require explicit consent from the identity owner before releasing information. While the economic implications of such systemic changes present significant implementation challenges, establishing this intermediary authentication layer would provide technically sophisticated users with substantially greater control over their personal credit information and enhanced protection against unauthorized access.